\documentclass[10pt,a4paper]{amsart}
\usepackage[margin=19mm]{geometry}
\usepackage{amsmath,amssymb,mathtools}
\usepackage[scr=rsfs,cal=euler]{mathalfa}
\usepackage{xfrac}
\usepackage[unicode,hidelinks,bookmarks=false]{hyperref}
\newcommand{\ie}{i.e.\ }
\newcommand{\cf}{cf.\ }
\newcommand{\resp}{respectively\ }
\newcommand{\Subsection}{\subsection}
\providecommand{\nobreakdash}{\nobreak\hbox{-}}
\setlength{\emergencystretch}{2em}
\multlinegap0pt
\usepackage{amssymb}
\usepackage{xcolor}
\usepackage{algorithm}\usepackage{algpseudocode}
\newtheorem{lemma}{Lemma}
\newcommand{\R}{\mathbb{R}}
\DeclareMathOperator{\graph}{graph}
\newcommand{\loc}{\mathrm{loc}}
\newcommand{\vp}{\varphi}
\title{On polynomial solutions to the minimal surface equation}
\author{Yifan Guo}
\date{Source: arXiv:2404.00115v2 (CC BY 4.0)}
\begin{document}
\maketitle

\noindent\emph{Source and licence.} On polynomial solutions to the minimal surface equation. Version arXiv:2404.00115v2 (CC BY 4.0), distributed by arXiv under the Creative Commons Attribution 4.0 International licence, http://creativecommons.org/licenses/by/4.0/, as recorded in the arXiv licence metadata for this exact version and shown on the abstract page https://arxiv.org/abs/2404.00115v2. Full licence notice: \texttt{licenses/arxiv-2404.00115-CC-BY-4.0.txt}.\par\smallskip

\noindent\textit{Editorial scope.} Counted unit: the article's lemma convgr (source0.tex lines 687--691). The article's complete original proof of it is delivered with it (source0.tex lines 692--740, one \texttt{proof} environment of 49 lines). No mathematics is written, repaired or re-derived: the statement and the proof are the article's own lines, in the article's own notation, and no retained line is substituted or re-flowed.  Retained uncounted as the source-local context the proof needs: lines 506--525; lines 663--667; lines 669--671.  Nothing was omitted from any retained span, so the package carries no omission record.  No retained line contains a citation, so the wrapper carries no bibliography and no bibliography entry.  No retained line contains a figure environment, an image-inclusion command or drawing code, so no image is delivered and the package declares no asset dependency.  Editorial formatting only, in addition: an AMS \texttt{amsart} preamble that restates the article's own declarations and macros used by the retained lines, namely the environments \texttt{lemma} on separate counters, together with the standard packages those lines need.  The retained environments print in this document as Lemma 1 in order of appearance, whereas the article prints the same items with the numbering of its own full document.  The complete native source of the article as distributed in the arXiv e-print for this exact version is bundled unchanged as \texttt{sources/156930/source0.tex}.
\begin{lemma}\label{expansion}
	Let $P=P_m+\cdots+P_1$ be a polynomial solution. Then 
	\begin{enumerate}
		
		\item[(i)] $P_m$ satisfies $LP_m=0$ where \[Lu:=|\nabla u|^2\Delta u-\sum_{i,j=1}^nu_{x_i}u_{x_j}u_{x_ix_j}.\]
		
		\item[(ii)] For any $t\neq 0$, the set $\{P_m=t\}$ is a smooth minimal hypersurface in $\R^n$ with $\nabla P_m$ nonvanishing on $\{P_m=t\}$.
		
		\item[(iii)] For any $t\neq 0$, the function $u=\frac{1}{|\nabla P_m|}$ is the velocity of the family of hypersurfaces $(\{P_m=t\})_{t\neq0}$ and hence $u$ is a Jacobi field on $\{P_m=t\}$ for $t\neq 0$ i.e. $\Delta_{\{P_m=t\}}u+|A_{\{P_m=t\}}|^2u=0$.
		
		\item[(iv)] $P_{m-1}$ satisfies 
		\begin{equation}
			\label{m-1}\begin{aligned}
				DL(P_m)P_{m-1}:=&|\nabla P_m|^2\Delta P_{m-1}+2\Delta P_m\nabla P_m\cdot \nabla P_{m-1}\\
				-\sum_{i,j=1}^n&\left( P_{m,x_i}P_{m,x_i}P_{m-1,x_ix_j}+2P_{m,x_j}P_{m,x_ix_j}P_{m-1,x_i}\right)=0.
			\end{aligned}
		\end{equation}
		\item[(v)] For any $t\neq 0$, the function $v=-\frac{P_{m-1}}{|\nabla P_m|}$ is a Jacobi field on $\{P_m=t\}$.
	\end{enumerate}
\end{lemma}

\begin{equation}\label{pt}
	\begin{aligned}
		G_t^{(\lambda)}(x)=&\lambda P(x/\lambda)-t\lambda^{-m+1}\\=&\lambda^{-m+1}(P_m(x)-t)+\lambda^{-m+2}P_{m-1}(x)+\cdots+P_1(x)
	\end{aligned}
\end{equation}

\begin{equation}\label{normal}
	\nu_{\graph G_t^{(\lambda)}}=\frac{(\nabla G_t^{(\lambda)}(x),-1)}{\sqrt{1+|\nabla G_t^{(\lambda)}(x)|^2}}.
\end{equation}

\begin{lemma}\label{convgr}
Let $P=P_m+\cdots+P_1$ be a polynomial. Let $t\neq 0$ and $G_t^{(\lambda)}$ be as in (\ref{pt}). Then \[\lim_{\lambda\to0}\eta_{\lambda\#}\graph P-t\lambda^{-m+1}e_{n+1}=\lim_{\lambda\to0}\left[\graph G_t^{(\lambda)}\right] = \partial[P_m<t]\times \R\] where the convergence is both in $I_{n}^\loc(\R^{n+1})$ as currents and locally in $C^1$ as normal graphs over $\{P_m=t\}\times\R$.

Moreover, the family of hypersurfaces $(\graph G_t^{(\lambda)})_{\lambda\in [0,\epsilon)}$ is smooth in $\lambda$ and the velocity of the family at $\lambda=0$ is $-\frac{P_{m-1}}{|\nabla P_m|}$ viewed as a function on $\{P_m=t\}\times\R$.
\end{lemma}

\begin{proof}

Since $t\neq0$, by Lemma \ref{expansion} (i), we define $\nu(z)=\frac{\nabla P_m(z)}{|\nabla P_m(z)|}$ for all $z\in \{P_m=t\}$. By the implicit function theorem, we can choose a neighborhood $U$ of $z_0$ so that the following are satisfied:

(i) We can write $U$ as $U=\psi(W\times (-\sigma,\sigma))$ where $\sigma>0$, $W=U\cap\{P_m=t\}$ and $\psi:W\times (-\sigma,\sigma)\to U$ is a diffeomorphism given by $\psi(z,r)=z+r\nu(z)$.

(ii) There exists $\delta>0$ so that for all $(z,r)\in W\times (-\sigma,\sigma)$ we have \[\nabla P_m(z+r\nu(z))\cdot \nu(z)\ge 2\delta >0.\]
Using (\ref{pt}), we compute for $\lambda$ sufficiently small
\[
\begin{aligned}
	\frac{d}{dr}G_t^{(\lambda)}(z+r\nu(z))=&\nabla G_t^{(\lambda)}(z+r\nu(z))\cdot \nu(z)\\
	\ge&\lambda^{-m+1}\nabla P_m(z+r\nu(z))\cdot \nu(z)-O(\lambda^{-m+2})\\
	\ge&\delta \lambda^{-m+1}>0.
\end{aligned}
\]
Thus $r\mapsto G_t^{(\lambda)}(z+r\nu(z))$ is increasing for all $-\sigma<r<\sigma$.

For $z\in W$, we have \[|G_t^{(\lambda)}(z)|=|\lambda^{-m+2}P_{m-1}(z)+\cdots +\lambda^{-1}P_2(z)+P_1(z)|\le C\lambda^{-m+2}.\]
Therefore, the interval $[-\sigma \delta \lambda^{-m+1}+C\lambda^{-m+2},\sigma \delta \lambda^{-m+1}-C\lambda^{-m+2}]$ is in the range of $G_t^{(\lambda)}$ on $U$ by the fundamental theorem of calculus. Note that these intervals are increasing and cover $\R$ as $\lambda\to0$. Let $R>0$ and $(z,y)\in W\times (-R,R)$ be fixed.  There exists $\lambda_0(R)>0$ so that for all $0<\lambda<\lambda_0$, $(-R,R)$ is in the range of $r\mapsto G_t^{(\lambda)}(z+r\nu(z))$ on $(-\sigma,\sigma)$. Since $r\mapsto G_t^{(\lambda)}(z+r\nu(z))$ is increasing, there is a unique $-\sigma<r_0<\sigma$ so that $G_t^{(\lambda)}(z+r_0\nu(z))=y$.

We write $r_0=\vp^{(\lambda)}(z,y)$ and $x=x(z,y)=z+\vp^{(\lambda)}(z,y)\nu(z)$. Then
\begin{align}\label{graph}
	(x,G_t^{(\lambda)}(x))=(z,y)+\vp^{(\lambda)}(z,y)(\nu(z),0).
\end{align}
This means that we can write $\left[\graph G_t^{(\lambda)}\right]\llcorner (U\times(-R,R))$ as the normal graph of $\vp^{(\lambda)}$ over $W \times (-R,R)$. Using (\ref{graph}), the chain rule and (\ref{pt}), we compute the gradient of $\vp^{(\lambda)}$ as follows 
\begin{align}
	\frac{\partial \vp^{(\lambda)}}{\partial y}(z,y)=&\frac{1}{\nabla G_t^{(\lambda)}(x)\cdot\nu(z)}=O(\lambda ^{m-1})\\
	\frac{\partial \vp^{(\lambda)}}{\partial z_j}(z,y)=&\sum_{i=1}^{n}\frac{-1}{\nabla G_t^{(\lambda)}(x)\cdot\nu(z)}\frac{\partial G_t^{(\lambda)}}{\partial x_i}(x)(e_{ij}+\vp^{(\lambda)}(z,y)\frac{\partial \nu^i}{\partial z_j}(z))\\
	=&\sum_{i=1}^{n}\frac{-1}{\nabla P_m(x)\cdot\nu(z)}\frac{\partial P_m}{\partial x_i}(x)(e_{ij}+\vp^{(\lambda)}(z,y)\frac{\partial \nu^i}{\partial z_j}(z))+O(\lambda)\label{dz}
\end{align}
where  $(e_{ij})_{i=1,\dots,n,j=1,\dots,n-1}$ is the projection matrix from $\R^n$ to $T_z\{P_m=t\}=\nu(z)^{\perp}$.

From (\ref{pt}), we have 
\begin{equation}\label{defining.vp}
	P_m(x)-t=\lambda^{m-1}y-\lambda P_{m-1}(x)-\cdots- \lambda^{m-1}P_1(x).
\end{equation} Hence  with $x=z+\vp^{(\lambda)}(z,y)\nu(z)$, we have $P_m(x)\to t$ as $\lambda\to0$. Since $\nabla P_m\neq0$ on $\{P_m=t\}$, we have $d(x,\{P_m=t\})\to 0$. Hence\[\lim_{\lambda\to 0}\vp^{(\lambda)}(z,y)= 0.\]
Since $(e_{ij})$ is the projection onto $T_z\{P_m=t\}$ we have $\sum_{i=1}^{n}e_{ij}\partial_{x_i}P_m(z)=0$. Thus the first term in (\ref{dz}) converges to 0 in view of $d(x,\{P_m=t\})\to 0$. Hence we see that $\graph G_t^{(\lambda)}$ converge to $(U\cap \{P_m=t\})\times\R$ as normal graphs in $C^1$. We observe that the unit normal of $\graph G_t^{(\lambda)}$ as in (\ref{normal}) converges to $\nu$ as $\lambda\to0$. Thus $\left[\graph G_t^{(\lambda)}\right]\to \partial[P_m<t]\times \R$ as currents.

To show that the velocity of the family $(\graph G_t^{(\lambda)})_{\lambda\in [0,\epsilon)}$ at $\lambda=0$ is $-\frac{P_{m-1}}{|\nabla P_m|}$, it suffices to compute $\frac{\partial \vp^{(\lambda)}}{\partial \lambda}|_{\lambda=0}$. By Taylor expansion around $z\in \{P_m=t\}$, we have
\[\begin{aligned}
	P_m(x)-t=&P_m(z+\vp^{(\lambda)}\nu(z))-P_m(z)\\
	=&\nabla P_m(z)\cdot\nu(z)\vp^{(\lambda)}+o(\vp^{(\lambda)})\\
	=&|\nabla P_m(z)|\vp^{(\lambda)}+o(\vp^{(\lambda)}).
\end{aligned}\]
Since right hand side of (\ref{defining.vp}) is of order $O(\lambda)$ and $|\nabla P_m(z)|\neq0$, we have $\vp^{(\lambda)}=O(\lambda)$.
We divide both sides of (\ref{defining.vp}) by $|\nabla P_m(z)|\lambda$ and obtain
\[\frac{\vp^{(\lambda)}}{\lambda}+o(1)=-\frac{P_{m-1}(x)}{|\nabla P_m(z)|}+O(\lambda).\]
Taking $\lambda\to0$ we get $\frac{\partial \vp^{(\lambda)}}{\partial \lambda}|_{\lambda=0}(z)=-\frac{P_{m-1}}{|\nabla P_m|}$. This finishes the proof.
\end{proof}

\end{document}
